\section{Introduction}
\label{sec:intro}

The monolithic scaling of deep neural networks has ushered in unprecedented empirical capabilities. However, training a single, massive model to generalize across all tasks is computationally prohibitive and prone to catastrophic forgetting. To circumvent this, modular parameter ensembling has emerged. By training parameter-efficient adapters (such as LoRAs) or task-specific vectors on isolated domains, researchers construct libraries of specialized expert models sharing a pre-trained backbone. At test-time, the challenge lies in dynamically ensembling these weights—often called \emph{test-time model merging}—to handle incoming heterogeneous input streams containing mixed-domain samples without manual task boundaries.

While early works relied on static uniform averaging, modern techniques predict sample-specific routing coefficients using auxiliary routing networks. However, from a rigorous statistical perspective, current ensembling and dynamic routing methods suffer from severe, foundational vulnerabilities:

\begin{figure}[t]
\centering
\includegraphics[width=\columnwidth]{fiosr_vs_baselines.png}
\caption{Joint model ensembling accuracy under varying stream batch sizes ($B=1$ to $512$). Parametric methods (Linear Unreg, QWS-Merge, and L3-Softmax) catastrophically collapse to near-uniform accuracy due to severe test-time overfitting, whereas the parameter-free frameworks (PFSR+MBH and our FIOSR) completely bypass optimization, maintaining superior ensembling stability across all batching regimes.}
\label{fig:stream_robustness}
\end{figure}

\begin{itemize}\setlength{\itemsep}{0.5pt}\setlength{\parskip}{0pt}
    \item \textbf{The Dynamic Routing Paradox and Few-Shot Overfitting:} Parametric routers—which project representation spaces using learned linear layers or wave superposition heuristics—are calibrated on tiny support sets (e.g., $N=64$ samples). Under such extreme data scarcity, these over-parameterized routers suffer from catastrophic overfitting, frequently assigning saturated, incorrect weights to test samples. Well-regularized models like L3-Softmax escape overfitting only by collapsing routing weights to non-adaptive uniform averages, negating the benefit of dynamic ensembling.
    
    \item \textbf{Vectorization Collapse under Sequential Streams:} Standard dynamic routers rely heavily on the implicit variance-smoothing effect of large batch averages. On single-sample sequential streams (batch size $B=1$), the absence of this smoothing mask causes predicted coefficients to fluctuate wildly. This phenomenon, which we term \emph{Vectorization Collapse}, results in severe ensembling accuracy drops, rendering such methods useless in practical, low-latency streaming applications.
    
    \item \textbf{Flat Euclidean Geometrical Misspecification:} Recent "parameter-free" subspace routing (PFSR) methods resolve the optimization bottleneck by projecting representations onto frozen expert class prototypes using standard cosine similarity. However, computing a raw Euclidean or cosine distance in the weight space implicitly assumes that the underlying parameter and representation manifolds are flat and isotropic. This is geometrically and information-theoretically misspecified. Fine-tuning on specialized tasks warps local representation manifolds, making different coordinates exhibit asymmetric sensitivities to noise and task-irrelevant variations.
\end{itemize}

To resolve these critical limitations, we propose \textbf{Fisher-Information Optimal Subspace Routing (FIOSR)}, a mathematically rigorous, training-free, and parameter-free dynamic ensembling framework. Rather than assuming a flat Euclidean weight space, we treat the neural network parameter space as a Riemannian manifold, where the local geometry is defined by the \emph{Fisher Information Metric}. The Fisher Information Matrix (FIM) is the unique, coordinate-invariant Riemannian metric on parameter spaces, representing the sensitivity of the model's predictive distribution to local perturbations.

By computing a smoothed and power-scaled diagonal empirical Fisher Information Matrix (dFIM) of the specialized expert models over a tiny calibration split, we construct an analytical local Riemannian metric tensor. When projecting test representations onto expert class prototypes, we replace standard cosine similarity with a \emph{Fisher-Weighted Cosine Similarity}. This formulation mathematically warps the projection coordinate space, naturally suppressing noisy or task-irrelevant dimensions (which have low Fisher values) while magnifying highly informative, discriminative coordinates.

To safeguard our framework against statistical biases and stream-level non-stationarity, we integrate two essential stabilization mechanics:
\begin{enumerate}\setlength{\itemsep}{0.5pt}\setlength{\parskip}{0pt}
    \item \textbf{Class-Size Scaling Calibration (CSC):} An analytical normalization of similarity coordinates that corrects for statistical maximum bias under highly asymmetric expert vocabulary dimensions.
    \item \textbf{Micro-Batch Homogenization (MBH):} An unsupervised batch-partitioning algorithm that groups incoming samples by their dominant Fisher-weighted task coordinates, shielding the merged model from \emph{heterogeneity collapse}.
\end{enumerate}

Through exhaustive statistical ensembling experiments conducted across 10 independent random seeds within a 192-dimensional synthetic Analytical Coordinate Sandbox, we demonstrate that FIOSR matches or exceeds the oracle ceilings of individual specialized experts ($\approx 100\%$ on MNIST and FashionMNIST, and near-ceiling stability on CIFAR-10 and SVHN). Crucially, because FIOSR requires zero trainable parameters and zero test-time optimization, it is completely immune to the Dynamic Routing Paradox and Vectorization Collapse. It maintains absolute stability and superior performance across all stream batch sizes ($B=1$ to $512$), outperforming complex parametric ensembling baselines by up to $40.7\%$.

In summary, our key contributions are:
\begin{itemize}\setlength{\itemsep}{0.5pt}\setlength{\parskip}{0pt}
    \item We introduce an information-geometric perspective to test-time model ensembling, formally analyzing weight and representation projections through the lens of Riemannian manifolds.
    \item We formulate a smoothed and power-scaled diagonal Fisher Information regularizer that serves as an analytical coordinate filter, suppressing task-irrelevant noise in parameter-free routing.
    \item We demonstrate that our training-free formulation, combined with CSC and MBH, is completely immune to overfitting and Vectorization Collapse under low-latency sequential streams.
    \item We provide a rigorous empirical evaluation against five major ensembling baselines across 10 random seeds, establishing the state-of-the-art for parameter-free test-time model merging.
\end{itemize}
